\begin{pastexam}[label={px:ntu-108-esoe-ds-02}]{臺大工科海洋所 108 學年度 第 2 題}
（程式設計；本卷規定程式題用 C++ 或 Java）

\textbf{(15 \%)} Write a method called \textsf{makeList} in a \textsf{SList} class with a tail. The
\textsf{makeList} function takes an array counts of \textsf{int}, constructs a singly-linked list,
and associates its head and tail pointers to the corresponding \textsf{SListNode} nodes. In the
constructed list, the first \textsf{counts[0]} items are the number \textsf{counts[0]}, the next
\textsf{counts[1]} items are the number \textsf{counts[1]}, etc. For example, if the input array is
\{1 3 2\}, then the output linked list is: \textsf{head}$\rightarrow$ \fbox{1} $\rightarrow$ \fbox{3}
$\rightarrow$ \fbox{3} $\rightarrow$ \fbox{3} $\rightarrow$ \fbox{2} $\rightarrow$ \fbox{2}
$\leftarrow$\textsf{tail}. Define the \textsf{SList} and \textsf{SListNode} classes to complete the
\textsf{makeList} method. When it is done, the \textsf{size} of nodes in \textsf{SList} should also
be updated accordingly.
\end{pastexam}

\begin{pastsoln}
卷首規定程式題用 C++ 或 Java，下面用 C++。 C++ 的陣列不帶長度，所以 \texttt{makeList} 多收一個
參數 \texttt{n}（\texttt{counts} 的長度）；用 Java 時就是 \texttt{counts.length}。 題目沒說呼叫時串列
原本是不是空的，這裡先把舊的節點清掉，再建新的串列。 \texttt{counts[i]} 是 $0$ 時，數字 \texttt{counts[i]}
不出現（迴圈一次都不跑）。

\begin{minted}[fontsize=\small]{cpp}
class SListNode {
public:
    int item;
    SListNode* next;
    SListNode(int i, SListNode* n = nullptr) : item(i), next(n) {}
};

class SList {
public:
    SListNode* head;
    SListNode* tail;
    int size;

    SList() : head(nullptr), tail(nullptr), size(0) {}
    ~SList() { clear(); }

    void clear() {                    // 釋放所有節點，變回空串列
        while (head != nullptr) {
            SListNode* p = head;
            head = head->next;
            delete p;
        }
        tail = nullptr;
        size = 0;
    }

    void makeList(const int counts[], int n) {
        clear();
        for (int i = 0; i < n; i++) {
            for (int j = 0; j < counts[i]; j++) {
                SListNode* node = new SListNode(counts[i]);
                if (tail == nullptr) {    // 第一個節點
                    head = node;
                } else {
                    tail->next = node;
                }
                tail = node;
                size++;
            }
        }
    }
};
\end{minted}

\textbf{正確性}。 有 \texttt{tail} 指標，新節點直接接在尾端（資結書第~5.2~節末的比較表：記了 \texttt{tail} 的單向串列，接在最後是 $O(1)$），不必每次
從頭走到尾。 內層迴圈的不變量：串列由前到後是 \texttt{counts[0]} 個 \texttt{counts[0]}、……、
\texttt{counts[i-1]} 個 \texttt{counts[i-1]}，再接 \texttt{j} 個 \texttt{counts[i]}；\texttt{head} 指第一個節點、
\texttt{tail} 指最後一個節點（串列空時兩者都是 \texttt{nullptr}），\texttt{size} 等於節點數。 接上一個節點時，
空串列要設 \texttt{head}，不空要讓舊的尾端指向它；兩種情形都把 \texttt{tail} 改成新節點、\texttt{size} 加一，
不變量保持。 新節點的 \texttt{next} 是 \texttt{nullptr}，所以串列的最後一個節點永遠指向 \texttt{nullptr}。
兩層迴圈跑完，串列就是題目要的樣子。

以 \texttt{\{1, 3, 2\}} 追蹤：$i = 0$ 時放一個 $1$，\texttt{head}、\texttt{tail} 都指向它；$i = 1$ 時接上三個
$3$；$i = 2$ 時接上兩個 $2$。 得到 $1 \to 3 \to 3 \to 3 \to 2 \to 2$，\texttt{tail} 指向最後的 $2$，
\texttt{size} $= 6$。

\textbf{複雜度}。 令 $N = \sum_i \texttt{counts[i]}$（建出來的節點數）。 每個節點 $O(1)$ 時間，外層迴圈
另外跑 $n$ 次，清掉舊串列和舊的節點數成正比，所以時間是 $\Theta(n + N)$ 加上清掉舊串列的時間；
額外空間就是 $N$ 個節點。 如果沒有 \texttt{tail}、每次從 \texttt{head} 走到尾端再接，時間會變成
$\Theta(N^2)$。
\end{pastsoln}
